% =====================================================================
% Section: Balancer Board Design and Droop Implementation Hardware
% Mock / planning-stage prose. Numeric values from CLAUDE.md, the IO CSV,
% and the BOM (rev 20260622); re-verify against the schematic PDF before
% publishing, especially anything flagged TODO(calibrate) in firmware.
% NOTE(orchestrator): bus nominal is quoted as 17.5 V below (rev-20260622
% design value). The 2026-07-08 plan / 2026-07-11 retune moved nominal to
% 16 V — reconcile with Section \ref{sec:droop} (which uses the 16 V-retune
% constants, V_0 = 15.91 V) before submission. One bus voltage, everywhere.
% =====================================================================

\section{Balancer Board Design and Droop Implementation Hardware}
\label{sec:board-design}

% NOTE(author): Open with the system-level framing: this is a purpose-built
% PCB (not an off-the-shelf BMS/charge-controller stack) because droop
% power sharing between two dissimilar sources needs an analog injection
% point that commercial boost modules don't expose. Argue that hardware
% co-design is a contribution in itself, not just a testbed.

\subsection{Overall Architecture}
\label{sec:board-architecture}

The balancer board (revision 20260622) implements a two-source,
bidirectional power-pathing architecture for a hybrid fuel-cell/battery
scale-model vehicle. Two independently regulated sources are combined onto
a common DC bus: a hydrogen fuel cell and a 2S lithium battery pack, each
boosted to the nominal bus voltage by an independent
TPS61288 synchronous boost converter (FC channel and BT channel). The two
boost outputs are not simply diode-ORed onto the bus; each is gated onto
the bus through a dedicated RT1987 ideal-diode power-path switch, so the
firmware -- not the analog hardware alone -- decides when each source is
allowed to contribute current.

% NOTE(author): Emphasize *why* ideal-diode switches instead of passive
% Schottky ORing: (1) negligible forward drop vs. a discrete diode at
% several amps, (2) true reverse-current blocking (no body-diode
% passthrough) so a disabled channel cannot be back-fed, which matters
% once regen braking is introduced, and (3) they are digitally
% controllable, which is the enabling feature for the sequencing state
% machine in Section~\ref{sec:power-pathing}.

Downstream of the bus, power flows through a further ideal-diode switch to
the VESC motor controller and the drive motor (\texttt{MOT\_PWR\_ENABLE}).
The motor is also the vehicle's regenerative braking source: during
deceleration the VESC can source current back onto the DC bus. This
current is intentionally \emph{not} routed back through the boost
converters -- doing so would require running a TPS61288 in reverse outside
its intended operating envelope and risks the back-feed hazard discussed
in Section~\ref{sec:power-pathing}. Instead, regenerated energy is routed
through a dedicated \texttt{REGEN\_ENABLE} path to the on-board battery
charger.

Battery charging is handled by two cooperating circuits with very
different bandwidths:

\begin{itemize}
  \item A \textbf{TL431/BSP170P shunt (chopper) clamp} is the
    \emph{fast}, board-local protection element. It is a fixed analog
    circuit (TL431 reference, BSP170P shunt MOSFET, dump resistor) with
    \emph{no firmware involvement} -- it exists to absorb the leading edge
    of a regen transient before any slower control loop can react, and to
    bound bus overvoltage independent of MCU state.
  \item A \textbf{Silvertel Ag105 MPPT module} is the \emph{secondary},
    slow harvester. It receives either fuel-cell bus power
    (\texttt{FC\_CHARGE\_ENABLE}) or regen power
    (\texttt{REGEN\_ENABLE}~+~\texttt{MOT\_PWR\_ENABLE}) and trickles it
    into the 2S pack under its own perturb-and-observe MPPT loop. Firmware
    controls it coarsely, through an active-low \texttt{MPPT\_DISABLE}
    GPIO and a two-register I\textsuperscript{2}C configuration write
    (voltage/current profile selection), not through a fast current-command
    interface.
\end{itemize}

% NOTE(author): This two-tier framing (fast passive clamp + slow active
% MPPT) is a good place to cite classic references on hybrid energy
% storage system power split (fast/slow filtering, e.g. supercapacitor
% + battery literature) even though the board doesn't literally have a
% supercapacitor -- the chopper is functionally the same kind of fast
% element.

A Teensy 4.1 (ARM Cortex-M7, \SI{600}{\mega\hertz}) is the sole controller
on the board. It executes the droop power-share control loop, the motor
current loop (talking to the VESC over UART), the sequencing state
machine for the ideal-diode switches, and the balancer/charger
supervisory logic, and exposes telemetry/command framing over UDP to a
Raspberry Pi bridge. A BQ29200 provides hardware-independent
overvoltage protection (OVP) on the individual battery cells; its
balance-enable pin is hardwired inactive, and only its OVP disable
(\texttt{CBAL\_DISABLE}) is under firmware control.

\subsection{Droop Implementation Hardware Chain}
\label{sec:droop-hardware}

% NOTE(author): This is the core novelty subsection -- spend the most
% care here. Argue explicitly that analog droop injection at the
% converter feedback node, commanded digitally by the MCU, is what lets
% a single control loop trade off "how much power-share authority do I
% have" against "how fast can I move it" without touching the boost
% converter's own fast voltage loop compensation.
% NOTE(orchestrator): the derivation/values live in Section \ref{sec:droop};
% keep this subsection qualitative (board-resource framing) to avoid
% duplication.

Each boost channel implements current-mode droop by perturbing its own
voltage-feedback divider in proportion to measured output current, so
that as a channel's current rises, its regulated bus voltage target sags
slightly -- the classical mechanism by which multiple sources sharing one
bus settle to a current split without a fast communication link between
them. On this board, the droop \emph{coefficient} is not fixed in
resistor value; it is a \emph{digitally commandable gain}, which is the
key hardware enabler for a closed-loop, MCU-mediated power-share
controller rather than a purely passive droop network. The signal chain
per channel is:

\begin{enumerate}
  \item \textbf{Current sense -- INA253A1IPWR.} A bidirectional,
    integrated-shunt current-sense amplifier sits at each boost channel's
    output, converting channel current to a ground-referenced analog
    voltage. The board specified the A3 variant (\SI{0.4}{\volt\per\ampere})
    but the A1 variant (\SI{0.1}{\volt\per\ampere}) was fitted by
    procurement error; firmware compensates with the corrected gain
    ($K_\text{sns} = \SI{0.1}{\volt\per\ampere}$).
  \item \textbf{Buffer/scale -- OPA197.} A rail-to-rail op-amp buffers and
    scales the sense voltage before it is (a) read back by the Teensy ADC
    for telemetry and share-error computation, and (b) used as the
    injection driver downstream of the MDAC.
  \item \textbf{Digitally-controlled attenuation -- AD5443 multiplying
    DAC.} A 12-bit MDAC (SPI, one per channel, independent chip-selects
    \texttt{CS\_MDAC\_FC} / \texttt{CS\_MDAC\_BT}) sits in the injection
    path and acts as a firmware-programmable attenuation
    element. This is what turns "droop coefficient" from a soldered
    resistor value into a run-time variable the power-share controller
    can command directly.
  \item \textbf{Feedback injection network.} The MDAC output is summed
    into each boost converter's own voltage-feedback divider through a
    dedicated injection resistor, so the converter's regulation setpoint
    is pulled down in proportion to (measured current) $\times$ (MDAC
    code), closing the droop loop at the analog level while the MCU sets
    only the gain.
\end{enumerate}

% NOTE(author): This is a good spot to state the board-level resource
% cost of droop as a design requirement, to motivate why this couldn't be
% bolted onto an off-the-shelf boost eval board: it needs (a) an isolated
% current-sense tap at each converter's output *before* the ideal-diode
% switch (so sensed current reflects only that channel's own contribution,
% not bus-shared current), (b) two independent SPI chip-selects plus a
% shared SPI bus (MOSI/MISO/SCK) reserved on the MCU, (c) two ADC channels
% per source (voltage division check + current sense) reserved out of the
% Teensy 4.1's analog pins, and (d) op-amp output headroom -- note that the
% OPA197 on this board is bodged to run from the 5 V rail rather than
% 3.3 V, which the droop-code-to-voltage mapping must account for.

Because the injection network taps into the converter's own compensated
feedback loop, the achievable droop range is bounded by the converter's
loop dynamics, not just by the MDAC's 12-bit resolution -- this is the
link between the board-level hardware described here and the plant model
used in the controller synthesis (Section~\ref{sec:robust}). The MCU-side
resource cost of this scheme is deliberately modest: one shared SPI bus,
two chip selects, two ADC channels per source for sense readback, and no
dedicated timer/PWM hardware -- the droop actuation is fully absorbed by
the analog injection network, leaving the MCU free to run the share
controller as a periodic software loop rather than a hardware peripheral.

\subsection{Power-Pathing: Ideal-Diode Switches and Sequencing}
\label{sec:power-pathing}

% NOTE(author): This subsection should read as a safety-case argument,
% not just a feature list -- the reviewer should come away understanding
% that the switch topology exists *because* naive "always-on" boost
% paralleling is unsafe with a bidirectional (regen) load, and that the
% sequencing rules are a direct, derivable consequence of that hazard,
% not arbitrary firmware policy.

Six RT1987 ideal-diode power-path switches, each independently controlled
by a Teensy GPIO, gate every power path on the board:

\begin{center}
\begin{tabular}{ll}
\texttt{FC\_BUS\_ENABLE} & FC boost $\rightarrow$ bus \\
\texttt{BT\_BUS\_ENABLE} & BT boost $\rightarrow$ bus \\
\texttt{MOT\_PWR\_ENABLE} & bus $\rightarrow$ VESC/motor node \\
\texttt{REGEN\_ENABLE} & regen node $\rightarrow$ charger \\
\texttt{FC\_CHARGE\_ENABLE} & bus (FC) $\rightarrow$ charger \\
\texttt{BT\_SEQUENCE\_ENABLE} & battery pack sequencing \\
\end{tabular}
\end{center}

This is a deliberate departure from a simpler "both boosts always
paralleled on the bus" design. Two hazards drive the sequencing
constraints:

\begin{itemize}
  \item \textbf{Back-feed / disabled-converter stress.} A converter that
    is disabled but still electrically connected to an energized bus can
    see reverse current or overshoot stress it was not designed for.
    Consequently, whenever the firmware brings a path up or down, the
    switch ordering is chosen so a regen or charge path is never left
    pointed at a converter that is not actively driving that node.
  \item \textbf{Hot-plug onto a discharged bulk capacitance.} Closing an
    ideal-diode switch onto a large discharged capacitor (the bus and,
    behind \texttt{MOT\_PWR\_ENABLE}, a \SI{470}{\micro\farad} node plus
    the VESC's own input capacitance) presents a near-short to the
    source side for the duration of the switch's soft-start ramp. If that
    ramp cannot supply the inrush within the switch's short-circuit
    protection window, the switch enters a burst-retry limit cycle, and
    the resulting repeated high-current transients have -- empirically,
    on this board -- destroyed boost converters (Section~\ref{sec:bringup}).
    This is why bus and motor-node bring-up is staged (switches first,
    allowed to settle, \emph{then} boosts; the motor node is pre-charged
    before \texttt{MOT\_PWR\_ENABLE} is closed) rather than done in one
    step.
\end{itemize}

The explicit rule set (mirrored in the firmware's power-path sequencing
guard) is: \texttt{BT\_SEQUENCE\_ENABLE} initializes off and is latched on
once during bring-up; \texttt{FC\_CHARGE\_ENABLE} requires
\texttt{BT\_BUS\_ENABLE} and \texttt{REGEN\_ENABLE} to both be low before
it is asserted (charger input path and regen path are mutually
exclusive); and \texttt{MOT\_PWR\_ENABLE} is only asserted onto a
pre-charged, sufficiently energized bus.

As a defense-in-depth measure independent of firmware correctness, every
enable net carries a \SI{10}{\kilo\ohm} pull-down to ground, so that
during MCU reset or any high-impedance boot window, every switch defaults
open (off) rather than floating into an undefined, possibly-closed state.

% NOTE(author): Good place to reference the incident log
% (docs/boost-bringup-debug.md) qualitatively -- "bench bring-up
% experience confirmed this hazard is not merely theoretical" -- the
% condensed version is Section \ref{sec:bringup}.

\subsection{Major Components}
\label{sec:bom-table}

% NOTE(author): Trim/reorder rows to taste; this table is meant to be a
% quick-reference companion, not a full BOM reproduction.

\begin{table}[H]
\centering
\caption{Major active components on the balancer board (rev.\ 20260622)
and their role in the droop / power-pathing architecture.}
\label{tab:major-ics}
\begin{adjustwidth}{-\extralength}{0cm}
\centering
\begin{tabular}{lll}
\toprule
\textbf{Component} & \textbf{Part} & \textbf{Role} \\
\midrule
FC / BT boost regulator & TI TPS61288 & Per-source boost to the DC bus, hosts droop injection \\
Ideal-diode power-path switch ($\times$6) & Richtek RT1987 & Digitally gated, reverse-blocking power path \\
Current-sense amplifier & TI INA253A1IPWR & Per-channel forward-current sense ($K_\text{sns}=0.1$\,V/A) \\
Droop MDAC & Analog Devices AD5443 & 12-bit, SPI-programmable droop-injection gain \\
MDAC output buffer & TI OPA197 & Rail-to-rail output driving injection network \\
Battery charger / MPPT & Silvertel Ag105 & Secondary harvester, I$^2$C config + \texttt{MPPT\_DISABLE} GPIO \\
Fast regen clamp & TI TL431 + Infineon BSP170P & Passive, board-local overvoltage / regen shunt \\
Cell OVP & TI BQ29200 & 2S cell overvoltage protection, \texttt{CBAL\_DISABLE} controlled \\
MCU & NXP i.MX RT1062 (Teensy 4.1) & Droop, sequencing, motor-loop, and telemetry controller \\
Motor controller & VESC (external) & UART motor current control, regen source \\
\bottomrule
\end{tabular}
\end{adjustwidth}
\end{table}

\subsection{Figures}
\label{sec:board-figures}

% NOTE(author): All three figures below are placeholders. Each TODO
% specifies exactly what the figure must communicate so a figure drawn
% later (in draw.io / KiCad / a photo) can be dropped in without
% re-deriving intent from scratch.

\begin{figure}[H]
  \centering
  % TODO: Draw a system block diagram showing:
  %   - Two source blocks: "Fuel Cell" and "2S Battery Pack"
  %   - FC boost (TPS61288) and BT boost (TPS61288) blocks downstream of
  %     each source, each annotated with its droop injection loop
  %     (INA253 -> OPA197 -> AD5443 -> feedback divider, drawn as a small
  %     closed loop back into the boost block, NOT as a separate top-level
  %     block, to show it's a local analog loop the MCU only sets the gain
  %     for)
  %   - FC_BUS_ENABLE / BT_BUS_ENABLE ideal-diode switches between each
  %     boost and a single "DC Bus" node
  %   - MOT_PWR_ENABLE switch from DC Bus to a "VESC + Motor" block
  %   - A regen arrow from "VESC + Motor" back toward the bus/regen node,
  %     clearly distinguished (dashed or colored) from the forward power
  %     flow arrows
  %   - REGEN_ENABLE and FC_CHARGE_ENABLE switches both feeding into the
  %     "Ag105 MPPT Charger" block, with a note that they are mutually
  %     exclusive
  %   - TL431/BSP170P "Fast Chopper Clamp" block shown directly on the
  %     regen node, NOT gated by any switch, to visually distinguish it
  %     from the firmware-controlled slow path
  %   - BQ29200 block on the battery pack with CBAL_DISABLE line from MCU
  %   - Single "Teensy 4.1" MCU block with lines fanning out to every
  %     enable/sense signal, labeled with the Code Names from the IO map
  \fbox{\parbox[c][5cm][c]{0.9\linewidth}{\centering PLACEHOLDER --- system
  block diagram; see TODO comment in source for required content}}
  \caption{Block diagram of the balancer board power and control
    architecture. \emph{[Placeholder -- diagram to be produced from the
    schematic.]}}
  \label{fig:block-diagram}
\end{figure}

\begin{figure}[H]
  \centering
  % TODO: Draw a power-path topology diagram (distinct from the block
  % diagram above -- this one should look like a simplified schematic /
  % single-line power diagram, not a signal-flow block diagram) showing:
  %   - Every RT1987 ideal-diode switch drawn as its actual topology
  %     symbol (or a diode-with-a-gate-control glyph), not a generic box
  %   - The six enable signals labeled at their respective switches using
  %     the exact Code Names
  %   - The 470uF bulk capacitor drawn explicitly at the motor/regen node
  %     to visually justify the hot-plug/pre-charge discussion
  %   - Hazard call-outs: (1) "back-feed risk" arrow at any
  %     switch-disabled-converter junction, (2) "inrush / hot-plug risk"
  %     bracket over the motor-node capacitance
  %   - A small inset legend distinguishing: always-open-by-default paths,
  %     latch-on-once paths (BT_SEQUENCE_ENABLE), and mutually-exclusive
  %     path pairs (FC_CHARGE_ENABLE vs. REGEN_ENABLE/BT_BUS_ENABLE)
  \fbox{\parbox[c][5cm][c]{0.9\linewidth}{\centering PLACEHOLDER ---
  power-path topology diagram; see TODO comment in source}}
  \caption{Power-path topology showing the six RT1987 ideal-diode
    switches, their sequencing constraints, and the hazard points that
    motivate them. \emph{[Placeholder -- to be produced from the
    schematic.]}}
  \label{fig:power-path-topology}
\end{figure}

\begin{figure}[H]
  \centering
  % TODO: Insert an annotated photo of the assembled board (rev 20260622)
  % once bench bring-up is far enough along to have a clean, populated
  % board shot (not mid-bodge). Annotate with call-out labels for:
  %   - FC and BT boost regulator sections
  %   - The RT1987 switch cluster
  %   - The Ag105 charger module footprint
  %   - The TL431/BSP170P chopper components
  %   - The BQ29200 balancer
  %   - The MDAC/OPA197 droop injection section
  %   - The Teensy 4.1 header location
  % If any bodge components (BT/FC output caps, RC-BT resistor swap,
  % RD1 215k) are visible/relevant, annotate them and note in the caption
  % that the photo reflects the as-bodged, not as-designed, board.
  \fbox{\parbox[c][5cm][c]{0.9\linewidth}{\centering PLACEHOLDER ---
  annotated board photo; see TODO comment in source}}
  \caption{Assembled balancer board (rev.\ 20260622), annotated.
    \emph{[Placeholder -- photo pending final bench assembly.]}}
  \label{fig:board-photo}
\end{figure}

% =====================================================================
% PLANNING NOTES
%
% 1. Numeric provenance: bus voltage, K_sns (0.1 V/A, INA253A1 variant
%    mixup), and the six enable-pin names are taken verbatim from
%    CLAUDE.md / the IO CSV / the BOM as of rev 20260622. Several firmware
%    constants referenced implicitly here (LIMIT_V_BUS_MAX, divider
%    scale factors, AG105_SETTLE_MS, BUS_SETTLE_MS) are still
%    TODO(calibrate) in the firmware itself -- do NOT state specific
%    values for those in the paper body until bench calibration lands.
%
% 2. BUS VOLTAGE INCONSISTENCY: this draft's original prose said 17.5 V;
%    the 2026-07-11 retune moved nominal to ~16 V (V_0 = 15.91 V in
%    Sec. \ref{sec:droop}). The explicit voltage has been removed from the
%    prose here; pick one number project-wide before submission.
%
% 3. The "hardware bodge" history (BT compensator resistor swap, output-
%    cap bodges, corrected failure analysis) is NOT in this section's
%    prose — it lives in Section \ref{sec:bringup}.
%
% 4. Unit macros: this section uses siunitx (\SI) and booktabs — both are
%    loaded via the main file preamble additions; verify they don't clash
%    with mdpi.cls.
%
% 5. Cross-check the component table part numbers/roles against the full
%    BOM CSV before submission (draft was built from BOM rows 1-102).
% =====================================================================
